% Motif search in GraphEasy (p1GraphEasy-con-AutoTuner) --- THEORY companion.
% Written 2026-09-24.  Every experiment referenced was executed on the tree.
\documentclass[11pt,a4paper]{article}

\usepackage[utf8]{inputenc}
\usepackage[T1]{fontenc}
\usepackage{lmodern}
\usepackage{amsmath,amssymb}
\usepackage{booktabs}
\usepackage{listings}
\usepackage{xcolor}
\usepackage[margin=1in]{geometry}
\usepackage[hidelinks]{hyperref}

\lstdefinestyle{plain}{basicstyle=\ttfamily\footnotesize,breaklines=true,frame=single,framesep=3pt,columns=fullflexible}

\newcommand{\code}[1]{\texttt{#1}}
\newcommand{\file}[1]{\texttt{#1}}
\newcommand{\Aut}{\operatorname{Aut}}

\title{\bfseries Motif Search in GraphEasy --- the Theory\\[2mm]
\large Backtracking, spanning orders, and symmetry breaking: how the implementation
relates to the textbook algorithms}
\author{GraphEasy (p1GraphEasy-con-AutoTuner)}
\date{September 2026}

\begin{document}
\maketitle

\begin{abstract}
\noindent
This is the theory companion to \file{MOTIF\_SEARCH\_EXPLAINED.tex}. It answers
one question in depth: \emph{what algorithm is actually running when you write
\code{a->b; b->c;} and ask for the count?} The short answer is a depth-first
backtracking search over the pattern's variables --- pick a vertex for \code{a},
then a vertex for \code{b}, check the required edges, recurse --- and the longer
answer is how the pattern's \emph{spanning structure} controls the enumeration,
how symmetry breaking turns counting into ``one representative per orbit'',
what the complexity is, and how this compares to VF2, RI, nauty-style canonical
labelling and ESU enumeration. Two verified experiments anchor the theory: a
$K_{2,2}$ bi-fan orbit of size 4 reported as exactly one match, and the effect
of variable order on the search plan.
\end{abstract}

\tableofcontents

%==============================================================================
\part{The mental model, checked against the code}

\section{``Loop over a, then loop over b, check the edge, recurse''}

The description in the question is essentially the algorithm. Concretely, for a
pattern with variables $a,b,c$:

\begin{enumerate}
  \item \textbf{Level 0 (variable \code{a}):} loop over candidate vertices.
  \item \textbf{Level 1 (variable \code{b}):} for each candidate, check the
        constraints that involve only \code{a} and \code{b} --- for
        \code{a->b} that means $\mathrm{kind}(b,a)$ and $\mathrm{kind}(a,b)$ must
        equal the pattern's requirements (both directions, so that ``no edge'' is
        also a check).
  \item \textbf{Level 2 (variable \code{c}):} same, against \code{a} and \code{b};
  \item \textbf{Recurse to depth $k$:} a complete assignment is a match; report it
        (or count it).
\end{enumerate}

That skeleton is exactly what \file{runtime.c:1271--1315} implements
(\code{motif\_collect\_backtrack\_runtime}) and what the emitted IR nest
implements (\file{MotifIRBuilder.cpp:246--293}). The one place where the two
backends differ is \emph{where the candidates come from}, and that is where the
``spanning tree'' in the question comes in.

\section{Where a spanning tree of the pattern enters --- and where it does not}

Think of the pattern as a graph whose vertices are the variables
$\{a,b,c,\dots\}$ and whose edges are the motif edges (plus, implicitly, the
``must be absent'' pairs, which are also constraints, just with the requirement
$0$). Two roles:

\begin{itemize}
  \item \textbf{Tree edges drive enumeration.} Once some variable is bound, a
        candidate for a variable adjacent to it in the pattern can be taken from
        that variable's adjacency list --- one CSR row --- instead of scanning
        all $n$ vertices. A connected order of the variables therefore gives
        every level after the first a row to walk. In the implementation, this is
        the \code{LevelPlan} (\file{MotifPattern.h}): \code{source = OutRow}
        with a \code{driverVar} (``candidate is an out-neighbour of
        $A[\text{prev}]$''), or \code{InRow} for the mirror case, or
        \code{FullScan} when no tree edge is available.
  \item \textbf{Non-tree (chord) edges are check-only constraints.} This is the
        standard pattern of subgraph-search algorithms: enumerate along a
        spanning tree, verify the remaining edges as constraints. In the
        implementation every pair $(j,\text{prev})$ appears in the level's
        \code{checks} list with \emph{both} directions; the driving edge is not
        removed from the list, because the row proves the edge \emph{exists} but
        not that its \emph{kind} (weight sign) matches.
\end{itemize}

What the implementation does \emph{not} do: it does not compute a minimum
spanning tree, nor any optimal order. The variable order is simply the order in
which the variables first appear in the motif edges (source before target,
\file{ASTBuilder.cpp:230--243}), and the driver for each level is chosen by a
recency heuristic: the most recently bound variable that has a required edge to
the current level (\file{MotifPattern.cpp:127--168}). Section 5 discusses what an
optimal order would be and why we can afford not to compute it.

\section{The two backends, in this language}

\begin{center}
\begin{tabular}{@{}p{0.30\linewidth}p{0.30\linewidth}p{0.30\linewidth}@{}}
\toprule
& runtime backend (default) & IR backend (\code{GRAPHEASY\_MOTIF\_BACKEND=ir}) \\
\midrule
candidate source & always $0..n-1$ & \code{FullScan}, or one CSR row per level
(\code{GRAPHEASY\_MOTIF\_ADJDRIVE=1}) \\
consistency test & \code{required} matrix built at run time & constants baked in,
same test, both directions \\
injectivity & \code{used[n]} array & $k$ pairwise compares \\
canonicality & $O(k!)$ search per complete assignment & lower-bound compares,
evaluated while descending \\
output & bindings (so \code{for each motif} works) & count only, today \\
\bottomrule
\end{tabular}
\end{center}

So the mental model ``loop over all vertices, then loop over b\,'' describes the
runtime backend verbatim, and describes the IR backend once adjacency driving is
enabled --- with the improvement that each level after the first walks a row
rather than the whole vertex set.

%==============================================================================
\part{The problem, formally}

\section{Definitions}

Let $G=(V,E)$ be a directed graph with a \emph{kind function}
\[
  \kappa_G(u,v)=
  \begin{cases}
    0 & \text{no edge } u\to v,\\
    1 & \text{edge with positive weight, or an unweighted edge},\\
    -1 & \text{edge with negative weight},\\
    2 & \text{edge with weight } 0.
  \end{cases}
\]

A \emph{pattern} $P$ has variables $X=\{x_0,\dots,x_{k-1}\}$ and a total
requirement matrix
\[
  R : X\times X \to \{0,+1,-1\},
  \qquad
  R(i,j)=
  \begin{cases}
    +1 & \text{if the motif wrote } x_i\to x_j,\\
    -1 & \text{if the motif wrote } x_i\mathrel{-\mid} x_j,\\
    0 & \text{otherwise (induced: the edge must be absent).}
  \end{cases}
\]

An \emph{embedding} is an injective map $\sigma:X\to V$ such that for all
$i\ne j$,
\[
  \kappa_G\bigl(\sigma(i),\sigma(j)\bigr)=R(i,j)
  \qquad\text{(and, when }R(i,j)=+1\text{ or }-1\text{, the first-hit edge weight has that sign).}
\]

The fourth line of the definition of $\kappa_G$ is worth reading twice: a
zero-weight edge satisfies \emph{neither} $+1$ nor $-1$ nor $0$ --- it poisons the
pair in both directions. This was verified experimentally
(\file{MOTIF\_SEARCH\_EXPLAINED.tex}, \S Stage 7): with a zero-weight chord, all
three candidate patterns return $0$).

\section{Counting modulo symmetry}

The number that a motif query reports is \emph{not} the number of injections
$\sigma$ satisfying the constraints. Let
\[
  \Aut(P)=\{\pi\in S_k : R(\pi(i),\pi(j)) = R(i,j)\ \forall i,j\}
\]
be the automorphism group of the pattern. Two embeddings $\sigma$ and
$\sigma\circ\pi$ with $\pi\in\Aut(P)$ use the same vertex set with permuted
variable names; a motif count reports one representative per orbit.

This is not pedantry --- it is what makes ``$K_{2,2}$ contains one bi-fan'' true
even though there are four injective bi-fan embeddings of a four-vertex
bipartite graph. Verified experiment (pattern
\code{a->x; a->y; b->x; b->y;} on the $K_{2,2}$ graph $0,1\to 2,3$):
\code{numMotifs} $=1$, and the reported binding was the canonical one.

%==============================================================================
\part{The search}

\section{Why depth-first backtracking is the right skeleton here}

The constraints of $R$ are \emph{binary} and \emph{monotone with respect to the
binding order}: the constraint between two variables is fully decided the moment
the later of the two is bound. Therefore a partial assignment that fails a
constraint can never be extended to a full embedding --- pruning is safe, and no
backtracking beyond the offending level is needed. This is the entire
justification for the ``check b; if b fails, next b'' structure.

The alternative enumeration strategy --- generate all $k$-vertex induced
subgraphs and classify each (ESU/Wernicke, Kavosh, MODA) --- is the right choice
when you want \emph{all} motifs of size $k$ at once. For a single fixed pattern,
direct backtracking examines far fewer partial states, because the pattern's
edges are used as filters from the first level on.

\section{Complexity}

Naively the search tree has $n$ branches per level and depth $k$, so
$O(n^k)$ leaves; each leaf costs $O(k^2)$ to check and $O(k!)$ to canonicalise in
the reference scheme. Three reductions, all implemented:

\begin{enumerate}
  \item \textbf{consistency pruning}: a level only expands candidates that
        satisfy all pairs to earlier variables; the branching factor at level $j$
        is the number of vertices admissible given the first $j$ choices;
  \item \textbf{adjacency driving}: with tree edges, the candidate set at level
        $j$ is one CSR row, so its size is the degree, not $n$; for sparse graphs
        this is the difference between $O(n^k)$ and $O\!\left(\prod_j d_j\right)$;
  \item \textbf{symmetry breaking}: the canonical lower bounds reject a candidate
        subtree as soon as it would produce a non-representative embedding, so
        the $|\Aut(P)|$-fold overcounting is removed \emph{during} the descent,
        not by a post-filter at the leaves.
\end{enumerate}

Worst case remains exponential in $k$ (subgraph isomorphism is NP-hard), but
$k\le 6$ is enforced by the semantic analyser
(\file{SemanticAnalyzer.cpp:987}), which is exactly the regime where a
pattern-specialised nest with constants and $k$ nested loops is practical.

\section{The canonicality reduction (proof of the fast form)}

\textbf{Reference test} (what the runtime does): an assignment $A$ is canonical
iff no automorphism image $A\circ\pi$ is lexicographically smaller than $A$,
where the comparison walks $i=0,1,\dots,k-1$.

\textbf{Fast test} (what the emitter does): for each non-identity
$\pi\in\Aut(P)$, let $p=\min\{i:\pi(i)\ne i\}$; require $A[p] < A[\pi(p)]$.

\begin{quote}
\emph{Claim.} For injective $A$, the two tests agree.

\emph{Proof.} For a given $\pi$, the permuted vector $B=A\circ\pi$ satisfies
$B[i]=A[\pi(i)]$. Since $A$ is injective and $\pi$ is a bijection,
$B[i]\ne A[i]$ exactly when $\pi(i)\ne i$; hence the first coordinate where $B$
can differ from $A$ is $p$. Lexicographic comparison therefore reduces to the
single comparison at $p$. Moreover $\pi(p)>p$: if $\pi(p)=q<p$ then, because
$p$ is the first moved index, $\pi(q)=q$; but then $\pi(q)=q=\pi(p)$ would
contradict injectivity of $\pi$. So $A[p]<A[\pi(p)]$ is a lower bound on the
coordinate $p$ --- which is bound \emph{before} $\pi(p)$ in the search --- and can
be evaluated while descending. \hfill$\square$
\end{quote}

\paragraph{Worked instance (verified).} The bi-fan pattern
\code{a->x; a->y; b->x; b->y;} has $|\Aut(P)|=4$: the sources $a,b$ are
interchangeable, the targets $x,y$ are interchangeable, independently. The
bound-producing automorphisms are $\pi_{ab}$ (first moved index $a$), $\pi_{xy}$
(first moved index $x$) and their product; the emitted nest therefore carries
exactly two lower bounds:
\[
  A[a] < A[b], \qquad A[x] < A[y].
\]
On the $K_{2,2}$ graph this drops the four embeddings to the single canonical
binding --- measured: \code{numMotifs} $=1$.

The FFL pattern \code{a->b; a->c; b->c;} has trivial automorphism group (the
directed triangle with a bypass has no symmetry: swapping $b,c$ would have to map
the present edge $b\to c$ to the absent edge $c\to b$), so no lower bounds are
emitted and both embeddings of the five-arc example are reported --- measured:
$2$, with bindings $(0,1,2)$ and $(0,1,3)$.

\section{Ordering: what a spanning tree does and does not buy}

\begin{itemize}
  \item A \emph{connected} variable order guarantees every level after the first
        has a row to walk (\code{OutRow}/\code{InRow}); a disconnected pattern
        forces a \code{FullScan} level, and --- under induced semantics --- all of
        that level's checks are ``must be non-adjacent'' constraints, which is
        expensive but correct.
  \item The \emph{choice of driver} within a level matters for cache behaviour
        (pick the most recently bound variable) and for the branching factor
        (pick a low-degree variable). The implementation picks by recency;
        VF2-style solvers additionally order variables by degree. At $k\le6$ one
        could brute-force all $k!\le720$ orders at compile time and keep the
        cheapest under a static cost model --- a natural, cheap future
        improvement, not a requirement today.
  \item The pattern's \emph{written} order matters for the plan: the variables are
        indexed in first-appearance order, so rewriting the motif edges changes
        the level plans (and speed). It does not change the answer --- verified:
        the FFL written as \code{b->c; a->b; a->c;} still counts $2$.
  \item A minimum spanning tree is not the right objective even in theory: for
        enumeration cost we would minimise the \emph{product of candidate-set
        sizes} (a treewidth-flavoured objective), not edge weights, and the
        pattern's edges carry no weights to minimise.
\end{itemize}

%==============================================================================
\part{Relation to the classical algorithms}

\begin{center}
\begin{tabular}{@{}p{0.22\linewidth}p{0.34\linewidth}p{0.36\linewidth}@{}}
\toprule
algorithm & idea & relation to this implementation \\
\midrule
plain backtracking & assign variables in order, check constraints & the runtime
backend, verbatim; the correctness reference \\
VF2 / VF2++ & feasibility sets, look-ahead, adaptive ordering & stronger pruning
and variable ordering; not implemented here --- at $k\le6$ the bookkeeping is
not repaid \\
RI (Bonnici et al.) & spanning tree + back-edge verification & exactly the
shape of the adjacency-driven IR nest: tree edges enumerate, chords check \\
nauty / McKay canonical labelling & canonically label the whole graph to dedupe
isomorphic subgraphs & we dedupe \emph{pattern orbits}, not graph isomorphism
classes; at $k\le6$ the orbit check is a handful of compares \\
ESU / Kavosh / MODA & enumerate all $k$-subgraphs, then classify & the better
tool for ``all motifs at once''; for one fixed pattern, direct matching wins \\
\bottomrule
\end{tabular}
\end{center}

The differential evidence that the pattern-level theory is implemented
consistently: on a 2000-vertex, 6000-arc graph the FFL count is $26$ under the
runtime backend, the IR nest, and the adjacency-driven nest --- three code paths,
one answer --- with per-run times $0.279$\,s, $0.219$\,s and $0.016$\,s
respectively (\file{MOTIF\_SEARCH\_EXPLAINED.tex}, \S Measured behaviour).

%==============================================================================
\part{Correctness obligations and the naming trap}

\section{Obligations}

\begin{itemize}
  \item \textbf{Soundness:} every reported/counted assignment satisfies all
        pair constraints in both directions, is injective, and is a canonical
        representative. Enforced by construction (the checks run before the
        recursion) and by the induced-edge test.
  \item \textbf{Completeness:} every canonical embedding is counted. The filters
        are \emph{necessary conditions} only, so no valid assignment is
        discarded; pruning removes only partial assignments that already violate
        a constraint.
  \item \textbf{Termination:} finite domains, finite depth, no cycles in the
        recursion.
  \item \textbf{Agreement of the two canonicality tests:} proved in \S 8; the
        reference test is kept in \file{MotifPattern} for unit-testing against a
        transcription of the runtime.
  \item \textbf{Agreement of the probe:} the emitted
        \code{\_\_ge\_motif\_edge\_kind} transcribes the runtime's
        \code{graph\_directed\_edge\_kind\_runtime}, including first-hit-wins and
        the zero-weight poison.
\end{itemize}

\section{The naming trap (verified)}

The \code{for each motif (...)} binding list is \textbf{positional}. The
variables of a pattern are collected in \emph{first-appearance order}, scanning
edges left to right and source before target (\file{ASTBuilder.cpp:230--243}).
Consequences:

\begin{itemize}
  \item for \code{\{ a->x; a->y; b->x; b->y; \}} the internal order is
        $(a,x,y,b)$ --- so the correct binding spelling is
        \code{for each motif (a,x,y,b)}, not \code{(a,b,x,y)};
  \item the semantic analyser only checks the \emph{arity}
        (\file{SemanticAnalyzer.cpp:687}), so a wrong order silently permutes the
        names.
\end{itemize}

\noindent\emph{Demonstration.} Pattern written as \code{\{ b->c; a->b; a->c; \}}
(internal order $b,c,a$) on the five-arc graph, printing
\code{hasEdge(G,a,b)} and \code{hasEdge(G,a,c)} per match, $2$ matches:

\begin{center}
\begin{tabular}{@{}lll@{}}
\toprule
binding spelling & printed pairs & meaning \\
\midrule
\code{(a,b,c)} (wrong order) & \code{1 0 / 1 0} & ``a'' is really $b$, ``c'' is really $a$; the second probe asks for an absent edge \\
\code{(b,c,a)} (first-appearance order) & \code{1 1 / 1 1} & both probes are pattern edges, as they should be \\
\bottomrule
\end{tabular}
\end{center}

\noindent Practical rule: write the motif edges so that the variables appear in
the same order as the \code{for each} list, or list them in first-appearance
order. (A friendlier compiler would match names instead of positions, or reject
a permutation; today it does neither.)

%==============================================================================
\appendix

\section{Annotated search trees}

\paragraph{FFL on the five-arc graph.} Pattern \code{a->b; a->c; b->c;},
$\Aut=1$, no bounds. Depth-first with levels $a,b,c$:

\begin{lstlisting}[style=plain]
a=0  accept (all pairs checked vacuously at level 0)
  b=1  kind(1,0)=0 ok, kind(0,1)=1 ok
    c=2  kind(2,0)=0 ok, kind(0,2)=1 ok, kind(2,1)=0 ok, kind(1,2)=1 ok  -> MATCH (0,1,2)
    c=3  ... same shape                                                -> MATCH (0,1,3)
    c=1,0 rejected by injectivity
  b=2  kind(2,0)=0 ok, kind(0,2)=1 ok
    c=3  kind(3,2)=0 ok, kind(2,3) must be 1 -> 0  -> reject  (no 2->3)
  b=3  kind(3,0)=0 ok, kind(0,3)=1 ok
    c=2  kind(2,3) must be 1 -> 0 -> reject
a=1  ...
a=2  ...
\end{lstlisting}

\paragraph{Bi-fan on $K_{2,2}$: the orbit of size four becomes one.} Pattern
\code{a->x; a->y; b->x; b->y;}, bounds $a<b$, $x<y$. The four embeddings
$(a,b,x,y)\in\{(0,1,2,3),(1,0,2,3),(0,1,3,2),(1,0,3,2)\}$ all satisfy the edge
constraints; the lower bounds accept only $(0,1,2,3)$. Measured count: $1$.

\section{Formula card}

\begin{center}
\begin{tabular}{@{}ll@{}}
\toprule
quantity & expression \\
\midrule
kind function & $\kappa_G(u,v)=0$ (no edge), $1$ (unweighted or $w>0$), $-1$ ($w<0$), $2$ ($w=0$) \\
embedding condition & $\kappa_G(\sigma(i),\sigma(j))=R(i,j)$ for all $i\ne j$, $\sigma$ injective \\
automorphism group & $\Aut(P)=\{\pi : R(\pi(i),\pi(j))=R(i,j)\}$ \\
canonical representative & $\sigma$ such that no $\sigma\circ\pi$ is lexicographically smaller \\
fast bounds & for each non-identity $\pi$: $A[p]<A[\pi(p)]$, $p=\min\{i:\pi(i)\ne i\}$ \\
brute-force cost & $O(n^k)$ leaves, $O(k^2)$ checks per leaf; row-driven levels replace $n$ by the degree \\
\bottomrule
\end{tabular}
\end{center}

\end{document}
